\subsection{Conic Sections and Conics}

\begin{tcolorbox}[title=Problem 1, breakable]
    Determine the equation of each of the circles with the following centre and radius.
\end{tcolorbox}
\begin{proof}[Solution]
    \textbf{(a)} Centre the origin, radius $1$:
    \[x^2 + y^2 = 1.\]

    \textbf{(b)} Centre the origin, radius $4$:
    \[x^2 + y^2 = 16.\]

    \textbf{(c)} Centre $(3,1)$, radius $2$:
    \[(x-3)^2 + (y-1)^2 = 4.\]

    \textbf{(d)} Centre $(3,4)$, radius $3$:
    \[(x-3)^2 + (y-4)^2 = 9.\]
\end{proof}

\begin{tcolorbox}[title=Problem 2, breakable]
    Determine the condition on the numbers $f$, $g$, and $h$ in the equation
    \[x^2 + y^2 + fx + gy + h = 0\]
    for the circle to pass through the origin.
\end{tcolorbox}
\begin{proof}[Solution]
    Clearly $(0,0)$ must be a solution, so we require
    \[0^2 + 0^2 + f(0) + g(0) + h = 0,\]
    which gives $h = 0$. Thus the condition is $h = 0$.
\end{proof}

\begin{tcolorbox}[title=Problem 3, breakable]
    Determine the centre and radius of the circles given by the following equations.
\end{tcolorbox}
\begin{proof}[Solution]
    \textbf{(a)} $x^2 + y^2 - 2x - 6y + 1 = 0$.
    By Theorem 2 we have centre
    \[\left(-\tfrac{1}{2}(-2),\, -\tfrac{1}{2}(-6)\right) = (1, 3)\]
    and radius
    \[\sqrt{\tfrac{1}{4}(-2)^2 + \tfrac{1}{4}(-6)^2 - 1} = \sqrt{1 + 9 - 1} = 3.\]

    \textbf{(b)} $3x^2 + 3y^2 - 12x - 48y = 0$. Dividing through by $3$ gives $x^2 + y^2 - 4x - 16y = 0$.
    By Theorem 2 we have centre
    \[\left(-\tfrac{1}{2}(-4),\, -\tfrac{1}{2}(-16)\right) = (2, 8)\]
    and radius
    \[\sqrt{\tfrac{1}{4}(-4)^2 + \tfrac{1}{4}(-16)^2} = \sqrt{4 + 64} = \sqrt{68} = 2\sqrt{17}.\]
\end{proof}

\begin{tcolorbox}[title=Problem 4, breakable]
    Determine the set of points $(x,y)$ in $\mathbb{R}^2$ that satisfies each of the following equations.
\end{tcolorbox}
\begin{proof}[Solution]
    \textbf{(a)} $x^2 + y^2 + x + y = 1$.
    By Theorem 2 we have centre $\left(-\tfrac{1}{2}, -\tfrac{1}{2}\right)$ and radius
    \[\sqrt{\tfrac{1}{4} + \tfrac{1}{4} + 1} = \sqrt{\tfrac{3}{2}}.\]
    So the set is all points $(x,y) \in \mathbb{R}^2$ with distance $\sqrt{\tfrac{3}{2}}$ from $\left(-\tfrac{1}{2}, -\tfrac{1}{2}\right)$, i.e.\ a circle.

    \textbf{(b)} $x^2 + y^2 - 2x + 4y + 5 = 0$.
    By Theorem 2 we have centre
    \[\left(-\tfrac{1}{2}(-2),\, -\tfrac{1}{2}(4)\right) = (1, -2)\]
    and radius
    \[\sqrt{\tfrac{1}{4}(-2)^2 + \tfrac{1}{4}(4)^2 - 5} = \sqrt{1 + 4 - 5} = 0.\]
    So the set is the single point $(1, -2)$.

    \textbf{(c)} $2x^2 + 2y^2 + x - 3y - 5 = 0$. Dividing through by $2$ gives
    \[x^2 + y^2 + \tfrac{1}{2}x - \tfrac{3}{2}y - \tfrac{5}{2} = 0.\]
    By Theorem 2 we have centre
    \[\left(-\tfrac{1}{2}\left(\tfrac{1}{2}\right),\, -\tfrac{1}{2}\left(-\tfrac{3}{2}\right)\right) = \left(-\tfrac{1}{4}, \tfrac{3}{4}\right)\]
    with radius
    \[\sqrt{\tfrac{1}{4}\left(\tfrac{1}{2}\right)^2 + \tfrac{1}{4}\left(\tfrac{3}{2}\right)^2 + \tfrac{5}{2}}
    = \sqrt{\tfrac{1}{16} + \tfrac{9}{16} + \tfrac{40}{16}} = \sqrt{\tfrac{50}{16}} = \tfrac{5\sqrt{2}}{4}.\]
    So the set is all points $(x,y) \in \mathbb{R}^2$ with distance $\tfrac{5\sqrt{2}}{4}$ from $\left(-\tfrac{1}{4}, \tfrac{3}{4}\right)$, i.e.\ a circle.
\end{proof}

\begin{tcolorbox}[title=Problem 5, breakable]
    Determine which, if any, of the following pairs of intersecting circles are mutually orthogonal.
\end{tcolorbox}
\begin{proof}[Solution]
    \textbf{(a)} $C_1 = \{(x,y) \mid x^2 + y^2 - 4x - 4y + 7 = 0\}$ and
    $C_2 = \{(x,y) \mid x^2 + y^2 + 2x - 8y + 5 = 0\}$.
    Notice
    \[(-4)(2) + (-4)(-8) = -8 + 32 = 24, \qquad 2(7 + 5) = 2(12) = 24.\]
    So by Theorem 3, $C_1$ and $C_2$ are orthogonal.

    \textbf{(b)} $C_1 = \{(x,y) \mid x^2 + y^2 + 3x - 6y + 5 = 0\}$ and
    $C_2 = \{(x,y) \mid 3x^2 + 3y^2 + 4x + y - 15 = 0\}$.
    Dividing through by $3$ on $C_2$ gives
    \[C_2 = \left\{(x,y) \mid x^2 + y^2 + \tfrac{4}{3}x + \tfrac{1}{3}y - 5\right\}.\]
    Notice
    \[(3)\!\left(\tfrac{4}{3}\right) + (-6)\!\left(\tfrac{1}{3}\right) = 4 - 2 = 2,
    \qquad 2\bigl(5 + (-5)\bigr) = 2(0) = 0.\]
    Since $2 \neq 0$, by Theorem 3, $C_1$ and $C_2$ are not orthogonal.
\end{proof}

\begin{tcolorbox}[title=Problem 6, breakable]
    Find the equation of the line through the points of intersection of the circles
    \[x^2 + y^2 - 3x + 4y - 1 = 0, \qquad x^2 + y^2 + 5x - 6y + 3 = 0.\]
\end{tcolorbox}
\begin{proof}[Solution]
    By Theorem 4 the equation of the line is the difference of the two circle equations:
    \[\bigl(x^2 + y^2 - 3x + 4y - 1\bigr) - \bigl(x^2 + y^2 + 5x - 6y + 3\bigr) = 0,\]
    which simplifies to
    \[-8x + 10y - 4 = 0, \qquad \text{i.e.} \qquad 4x - 5y + 2 = 0.\]
\end{proof}

\begin{tcolorbox}[title=Problem 7, breakable]
    This question concerns the parabola $E$ with equation $y^2 = x$ and parametric equations
    $x = t^2$, $y = t$ where $t \in \mathbb{R}$.

    \textbf{(a)} Write down the focus, vertex, axis, and directrix of $E$.
\end{tcolorbox}
\begin{proof}[Solution]
    Writing $y^2 = 4ax = x$, so $4a = 1 \implies a = \tfrac{1}{4}$.
    So it has focus $\left(\tfrac{1}{4}, 0\right)$ and directrix $x = -\tfrac{1}{4}$.
    Its vertex is $(0,0)$ and its axis is the $x$-axis.

    \medskip
    Let $P$ and $Q$ be coordinates on $E$ so that
    \[P = (t_1^2, t_1), \qquad Q = (t_2^2, t_2).\]
    If $t_1 \neq t_2$ then the chord $PQ$ has slope
    \[\frac{t_2 - t_1}{t_2^2 - t_1^2} = \frac{1}{t_2 + t_1}.\]
    Since $(t_1^2, t_1)$ lies on the line we have
    \[(y - t_1) = \left(\frac{t_2 - t_1}{t_2^2 - t_1^2}\right)(x - t_1^2)
    \implies (y - t_1)(t_2^2 - t_1^2) = (t_2 - t_1)(x - t_1^2)
    \implies (y - t_1)(t_2 + t_1) = (x - t_1^2).\]
    It passes through the focus $\left(\tfrac{1}{4}, 0\right)$ iff
    \[(0 - t_1)(t_2 + t_1) = \left(\tfrac{1}{4} - t_1^2\right)
    \implies -t_1 t_2 - t_1^2 = \tfrac{1}{4} - t_1^2
    \implies -t_1 t_2 = \tfrac{1}{4}
    \implies t_1 t_2 = -\tfrac{1}{4}.\]
    So the chord $PQ$ is a focal chord precisely when $t_1 t_2 = -\tfrac{1}{4}$.
\end{proof}

\begin{tcolorbox}[title=Problem 8, breakable]
    Let $P$ be an arbitrary point on the ellipse with equation
    \[\frac{x^2}{a^2} + \frac{y^2}{b^2} = 1\]
    and focus $F(ae, 0)$. Let $M$ be the midpoint of $FP$. Prove that $M$ lies on an
    ellipse whose center is midway between the origin and $F$.
\end{tcolorbox}
\begin{proof}[Solution]
    We now prove $M$ lies on the ellipse
    \[\frac{\left(x - \tfrac{ae}{2}\right)^2}{\left(\tfrac{a}{2}\right)^2}
    + \frac{y^2}{\left(\tfrac{b}{2}\right)^2} = 1,\]
    which is centered at $\left(\tfrac{ae}{2}, 0\right)$, the midpoint of the origin and $F$.
    Parametrize the original ellipse by $P = (a\cos t,\, b\sin t)$ with $t \in \mathbb{R}$,
    so every point $P$ on it has this form. With $F = (ae, 0)$, the midpoint of $FP$ is
    \[M = \left(\frac{ae + a\cos t}{2},\ \frac{b\sin t}{2}\right).\]
    Substituting the coordinates of $M$ into the left-hand side, we have
    \[x - \frac{ae}{2} = \frac{ae + a\cos t}{2} - \frac{ae}{2} = \frac{a\cos t}{2},\]
    so
    \[\frac{\left(x - \tfrac{ae}{2}\right)^2}{\left(\tfrac{a}{2}\right)^2}
    = \frac{\left(\tfrac{a\cos t}{2}\right)^2}{\left(\tfrac{a}{2}\right)^2} = \cos^2 t,
    \qquad
    \frac{y^2}{\left(\tfrac{b}{2}\right)^2}
    = \frac{\left(\tfrac{b\sin t}{2}\right)^2}{\left(\tfrac{b}{2}\right)^2} = \sin^2 t.\]
    Adding these gives
    \[\frac{\left(x - \tfrac{ae}{2}\right)^2}{\left(\tfrac{a}{2}\right)^2}
    + \frac{y^2}{\left(\tfrac{b}{2}\right)^2} = \cos^2 t + \sin^2 t = 1.\]
\end{proof}

\newpage
\begin{tcolorbox}[title=Problem 9, breakable]
    Let $P$ be a point $\left(\sec t,\ \tfrac{1}{\sqrt{2}}\tan t\right)$, where
    $t \in \left(-\tfrac{\pi}{2}, \tfrac{\pi}{2}\right) \cup \left(\tfrac{\pi}{2}, \tfrac{3\pi}{2}\right)$,
    on the hyperbola $E$ with equation $x^2 - 2y^2 = 1$.
    \\[2pt]
    \textbf{(a)} Determine the foci $F$ and $F'$ of $E$.\\
    \textbf{(b)} Determine the gradients $FP$ and $F'P$, when these lines are not parallel to the $y$-axis.\\
    \textbf{(c)} Determine the point $P$ in the first quadrant on $E$ for which $FP$ is perpendicular to $F'P$.
\end{tcolorbox}

\begin{proof}
    We compare $x^2 - 2y^2 = 1$ with
    $\dfrac{x^2}{a^2} - \dfrac{y^2}{b^2} = 1$. The vertices occur at $y = 0$, where
    $x^2 = 1$, so $a^2 = 1$ and $a = 1$. Matching the $y$-term,
    \[2y^2 = \frac{y^2}{b^2} \implies b^2 = \frac{1}{2}.\]
    For a hyperbola $b^2 = a^2(e^2 - 1)$, so
    \[\frac{1}{2} = 1\cdot(e^2 - 1) \implies e^2 = \frac{3}{2}
    \implies e = \sqrt{\tfrac{3}{2}} = \frac{\sqrt{6}}{2}.\]
    The foci are at $(\pm ae, 0) = \left(\pm\dfrac{\sqrt{6}}{2},\, 0\right)$. Write
    \[F = \left(\frac{\sqrt{6}}{2}, 0\right), \qquad F' = \left(-\frac{\sqrt{6}}{2}, 0\right).\]
\end{proof}

\begin{proof}
    With $P = \left(\sec t,\ \tfrac{1}{\sqrt 2}\tan t\right)$, the gradient of $FP$ is
    \[m_{FP} = \frac{\tfrac{1}{\sqrt{2}}\tan t - 0}{\sec t - \tfrac{\sqrt{6}}{2}}
    = \frac{\tfrac{1}{\sqrt{2}}\tan t}{\sec t - \tfrac{\sqrt{6}}{2}},\]
    and the gradient of $F'P$ is
    \[m_{F'P} = \frac{\tfrac{1}{\sqrt{2}}\tan t - 0}{\sec t + \tfrac{\sqrt{6}}{2}}
    = \frac{\tfrac{1}{\sqrt{2}}\tan t}{\sec t + \tfrac{\sqrt{6}}{2}}.\]
\end{proof}

\newpage
\begin{proof}
    The lines $FP$ and $F'P$ are perpendicular iff the vectors $\vec{FP}$ and $\vec{F'P}$
    With
    $P = \left(\sec t,\ \tfrac{1}{\sqrt{2}}\tan t\right)$,
    $F = \left(\tfrac{\sqrt{6}}{2}, 0\right)$, and $F' = \left(-\tfrac{\sqrt{6}}{2}, 0\right)$,
    \[\vec{FP} = \left(\sec t - \frac{\sqrt{6}}{2},\ \frac{1}{\sqrt{2}}\tan t\right),
    \qquad
    \vec{F'P} = \left(\sec t + \frac{\sqrt{6}}{2},\ \frac{1}{\sqrt{2}}\tan t\right).\]
    Their dot product is
    \[\vec{FP} \cdot \vec{F'P}
    = \left(\sec t - \frac{\sqrt{6}}{2}\right)\left(\sec t + \frac{\sqrt{6}}{2}\right)
    + \frac{1}{2}\tan^2 t
    = \sec^2 t - \frac{3}{2} + \frac{1}{2}\tan^2 t.\]
    Setting this equal to $0$ and using $\tan^2 t = \sec^2 t - 1$,
    \[\sec^2 t - \frac{3}{2} + \frac{1}{2}\left(\sec^2 t - 1\right) = 0
    \implies \frac{3}{2}\sec^2 t - 2 = 0
    \implies \frac{3}{2}\sec^2 t = 2,\]
    so $\sec^2 t = \tfrac{4}{3}$. In the first quadrant $\sec t > 0$, giving
    \[x = \sec t = \frac{2}{\sqrt{3}}, \qquad \tan^2 t = \sec^2 t - 1 = \frac{1}{3},\]
    so with $\tan t > 0$,
    \[y = \frac{1}{\sqrt{2}}\tan t = \frac{1}{\sqrt{2}}\cdot\frac{1}{\sqrt{3}} = \frac{1}{\sqrt{6}}.\]
    Thus
    \[P = \left(\frac{2}{\sqrt{3}},\ \frac{1}{\sqrt{6}}\right).\]
\end{proof}

\subsection{Properties of Conics}

\begin{tcolorbox}[title=Problem 1, breakable]
    Determine the gradient of the tangent to the curve in $\mathbb{R}^2$
    with parametric equations 
    \[x = 2 \cos t + \cos 2t + 1, y = 2 \sin t + \sin 2t\]
    at the point with paramter $t$, where $t$ is not a multiple of $pi$
    Hence determine the equation of the tangent to this curve at the points 
    with parameter $t = \pi/3$ and $t = \pi/2$.
\end{tcolorbox}

\begin{proof}
    By Theorem 1, the gradient of the tangent is 
    \[\frac{y'(t)}{x'(t)} = \frac{2\cos t + 2\cos(2t)}{-2\sin t - 2 \sin 2t}
    = -\frac{\cos t + \cos 2t}{\sin t + \sin 2t}.\]
    Then plugging in $t = \pi/3$ we have, since $\cos(\pi/3)+\cos(2\pi/3) = \tfrac12 - \tfrac12 = 0$,
    that the gradient is $0$, while $x = 3/2$ and $y = 3\sqrt{3}/2$, so the tangent is 
    \[y = \frac{3\sqrt{3}}{2}.\]   
    and plugging in $t = \pi/2$ we have, since $\cos(\pi/2)+\cos(\pi) = 0 - 1 = -1$ and $\sin(\pi/2)+\sin(\pi) = 1 + 0 = 1$, that the gradient is $1$, while $x = 0$ and $y = 2$, so the tangent is 
    \[y = x + 2.\]
\end{proof}

\begin{tcolorbox}[title=Problem 2, breakable]
    \begin{enumerate}
        \item Determine the equation of the tangent at a point $P$ on the rectangular parabola with 
        paramteric equations $x = t, y = 1/t$.
        \item Hence determine the equations of the two tangents to the rectangular hyperbola from the point $(1, -1)$.
    \end{enumerate}
\end{tcolorbox}